Let $R$ be a ring and $M$ an $R$-module. A basis of $M$ is a set $\{x_\alpha\}$ of elements of $M$ such that no proper submodule of $M$ contains all the $x_\alpha$. The module $M$ is said to have a finite basis, or to be a finite $R$-module, if it has a basis consisting of a finite number of elements. It is said to be cyclic if it has a basis consisting of a single element.

If $\{x_\alpha\}$ is any set of elements of $M$, then the smallest submodule of $M$ which contains all the $x_\alpha$ consists of those elements of $M$ which can be written in the form of a finite sum
$$a_1 x_{\alpha_1} + a_2 x_{\alpha_2} + \cdots + a_n x_{\alpha_n} + m_1 x_{\alpha_1} + m_2 x_{\alpha_2} + \cdots + m_n x_{\alpha_n}$$
where $a_i \in R$ and $m_i$ are integers.
